% ==============================================================================
% T0 Theory: Document 167
% Title:  LiNbO3 Material Properties on the xi-Power Ladder —
%         Phase Matching in PPLN: Testing a Geometric Assignment
% Author: Johann Pascher
% Date:   March 2026
% Ref:    Documents 041, 056, 061, 081, 122, 133, 161, 162, 165, 166
% ==============================================================================

\section*{Abstract}

Periodically poled lithium niobate (PPLN) is the physical core of the
Coherent Ising Machine (CIM): the $\chi^{(2)}$ nonlinear crystal implements,
via parametric amplification, the DOPO spin — the fundamental bit of the
Ising computer. The poling period $\Lambda$, refractive index $n(\omega)$,
and electro-optic coefficient $r_{33}$ are conventionally regarded as
empirically determined material parameters with no geometric origin.

This document tests whether a central parameter of LiNbO$_3$, the dispersion
difference in DOPO operation, lies on the $\xipar$-power ladder of T0 theory,
namely in the form:

\begin{equation}
    \Delta n_{(1064 \to 2128\,\text{nm})} \;\stackrel{?}{=}\; \sqrt{3}\cdot\xipar^{1/2} = 0.020
    \label{eq:main}
\end{equation}

with $\xipar = 4/30\,000$. This assignment does not hold: the Sellmeier equation
after Zelmon 1997 (congruent LiNbO$_3$, extraordinary) gives
$n_e(1064\,\text{nm}) = 2.1555$ and $n_e(2128\,\text{nm}) = 2.1219$, i.e.\
$\Delta n = 0.0336$. Only this value matches the poling period:
$\lambda_p/\Delta n = 1.064\,\mu\text{m}/0.0336 = 31.6\,\mu$m ($\Lambda = 31\,\mu$m),
whereas $0.020$ would require $\Lambda = 53\,\mu$m. $\sqrt{3}\,\xipar^{1/2} = 0.020$
therefore lies $41\,\%$ below the measured value. The exact equality
$\sqrt{3}\cdot(4/30000)^{1/2} = \sqrt{12/30000} = 0.02$ is an algebraic identity;
the measured value is known to only 2--3 digits.

In T0 theory the exponent $1/2$ appears as the weak coupling constant
$\alpha_W = \xipar^{1/2}$, as the lepton-mass ratio $m_e/m_\mu \propto \xipar^{1/2}$
and in further contexts (Docs.~041, 056, 061, 122, 133). The relations
$E_\mathrm{EW} \sim E_P \cdot \xipar^{1/2}$ and $m_h/M_P = \xipar^{1/2}$ are not
tenable. The PPLN dispersion does not lie on this rung; what remains as a match on
the $\xipar$ ladder is the ratio of the electro-optic coefficients
$r_{33}/d_{33} \approx 1 + \xipar^{1/4}$ (deviation $2.2\,\%$).

\clearpage

% ==============================================================================
\section{Introduction: PPLN as the Core of the Coherent Ising Machine}
\label{sec:intro}
% ==============================================================================

\subsection{The DOPO Mechanism}

The Coherent Ising Machine (CIM) solves the Ising Hamiltonian
$H = -\sum_{ij} J_{ij}\sigma_i\sigma_j$ via a network of DOPO pulses
(Degenerate Optical Parametric Oscillators) inside a fibre cavity.

The physical core is the $\chi^{(2)}$ nonlinear crystal: a pump photon at
frequency $\omega_p$ is converted into two signal and idler photons at
$\omega_p/2$. Above threshold the DOPO develops exactly two stable states
— phase 0 or $\pi$ — corresponding to the Ising spin $\sigma_i = \pm 1$.

The material of choice is periodically poled lithium niobate (PPLN). The
question of why LiNbO$_3$ is usually answered empirically: largest known
$d_{33}$ coefficient combined with low absorption losses. This document
tests whether a geometric answer follows from the $\xipar$-power ladder.

\subsection{The Phase-Matching Condition}

In the quasi-phase-matching (QPM) scheme the phase mismatch
$\Delta k = k_p - k_s - k_i \neq 0$ is compensated by periodically
inverting the crystal domains with period $\Lambda$:

\begin{equation}
    \Lambda = \frac{2\pi}{\Delta k} = \frac{\lambda_{\text{pump}}}{\Delta n}
    \label{eq:poling}
\end{equation}

The poling period is directly proportional to the inverse dispersion
difference $\Delta n = n(\lambda_{\text{pump}}) - n(2\lambda_{\text{pump}})$ ($\lambda_{\text{pump}}$: pump wavelength for the DOPO, SH wavelength for SHG).
For the DOPO, $\Delta k = 2\pi(n_p - n_s)/\lambda_p$, hence $\Lambda = \lambda_p/(n_p - n_s)$.

\begin{table}[H]
    \centering
    \adjustbox{max width=\textwidth}{%
\begin{tabular}{lccc}
        \toprule
        \textbf{Process} & \textbf{Pump} &
        \textbf{$\Lambda$\,($\mu$m)} & \textbf{$\Delta n$} \\
        \midrule
        OPO — DOPO operation (CIM)  & $1064\,\text{nm}$ & $31.0$ & $0.034$ \\
        SHG — Telecommunications    & $780\,\text{nm}$  & $19.5$ & $0.040$ \\
        SHG — Green laser           & $532\,\text{nm}$  & $6.8$ & $0.078$ \\
        \bottomrule
    \end{tabular}%
}
    \caption{Typical PPLN parameters for LiNbO$_3$ (extraordinary polarisation,
    room temperature, Sellmeier coefficients after Zelmon 1997).}
    \label{tab:ppln}
\end{table}
In detail, Sellmeier (Zelmon 1997, $n_e$) gives: DOPO $1064 \to 2128$~nm
$\Delta n = 0.0336$, $\Lambda = 31.6\,\mu$m; SHG $1560 \to 780$~nm $\Delta n = 0.0405$,
$\Lambda = 19.3\,\mu$m; SHG $1064 \to 532$~nm $\Delta n = 0.078$,
$\Lambda = 0.532/0.078 = 6.8\,\mu$m.

\clearpage

% ==============================================================================
\section{The Exponent \texorpdfstring{$\xipar^{1/2}$}{xi\^{}(1/2)} in T0 Theory}
\label{sec:xi_half}
% ==============================================================================

\subsection{%
  \texorpdfstring{$\xipar^{1/2}$}{xi\^{}(1/2)} as the Fundamental
  Electroweak Exponent}

In T0 theory, $\xipar^{1/2} = (4/30\,000)^{1/2} = 0.011547\ldots$ is not
merely one among many $\xipar$ powers but the fundamental transition exponent
between the Planck scale and particle physics. It appears in a remarkable
variety of independent contexts:

\begin{table}[H]
    \centering
    \adjustbox{max width=\textwidth}{%
\begin{tabular}{llll}
        \toprule
        \textbf{Physical quantity} & \textbf{T0 expression} &
        \textbf{Value} & \textbf{Doc.} \\
        \midrule
        Weak coupling constant &
            $\alpha_W = \xipar^{1/2}$ &
            $1.15 \times 10^{-2}$ & 041, 061 \\
        Lepton-mass ratio &
            $m_e/m_\mu \propto \xipar^{1/2}$ &
            $4.84 \times 10^{-3}$ & 056, 122 \\
        Bottom-quark mass &
            $m_b = v \cdot \tfrac{3}{2}\cdot\xipar^{1/2}$ &
            $4.25\,\text{GeV}$ & 041 \\
        Planck constant (scaling) &
            $\hbar \sim \sqrt{\xipar}$ & — & 105 \\
        \bottomrule
    \end{tabular}%
}
    \caption{The exponent $\xipar^{1/2} = 0.011547$ as the universal
    electroweak exponent in T0 theory. It marks the geometric transition
    from the Planck world to the particle-physics scale.}
    \label{tab:xi12}
\end{table}

Two further $\xipar^{1/2}$ relations do not hold: $m_h/M_P = \xipar^{1/2}$ gives
$1.15\times10^{-2}$, whereas the measured value is $m_h/M_P = 1.0\times10^{-17}$ (see
Doc.~068); $E_\mathrm{EW} \sim E_P\cdot\xipar^{1/2}$ gives about $1.4\times10^{17}$~GeV
with $E_P = 1.22\times10^{19}$~GeV, not $\sim 100$~GeV, and is therefore not tenable as
a derivation of the electroweak scale.

\subsection{Physical Meaning of the Exponent}

The exponent $1/2$ marks the \emph{electroweak transition} on the
$\xipar$-power ladder:

\begin{align}
    \text{Planck scale:}\quad          & \xipar^0     = 1
        && E \sim E_P \notag \\
    \text{Electroweak transition:}\quad & \xipar^{1/2} \approx 10^{-2}
        && \alpha_W,\; m_e/m_\mu
        \label{eq:rung} \\
    \text{Particle scale:}\quad        & \xipar^1     \approx 10^{-4}
        && m_e,\; \alpha_\mathrm{EM} \notag
\end{align}

The PPLN dispersion does not lie on this rung (Section~\ref{sec:main_result}).

\subsection{Distinction from Document 166}

Document~166 uses the exponent $1/4$ for the Casimir scale
($L_\xi \propto \xipar^{1/4}$, decomposition $41/4 = 10 + 1/4$), not $\xipar^{1/2}$;
there is no common exponent of the Casimir scale and the PPLN dispersion.

\clearpage

% ==============================================================================
\section{Test: \texorpdfstring{$\Delta n$ versus $\sqrt{3}\cdot\xipar^{1/2}$}{%
  Delta n vs. sqrt(3) xi\^{}(1/2)}}
\label{sec:main_result}
% ==============================================================================

\subsection{Numerical Comparison}

The dispersion difference of LiNbO$_3$ for the DOPO operation is obtained
from the Sellmeier refractive indices for the extraordinary polarisation
(Zelmon 1997):

\begin{equation}
    \Delta n = n_e(1064\,\text{nm}) - n_e(2128\,\text{nm})
    = 2.1555 - 2.1219 = 0.0336
\end{equation}

The expression of Eq.~\eqref{eq:main} gives:

\begin{align}
    \sqrt{3}\cdot\xipar^{1/2}
    &= \sqrt{12/30000} = 0.020000\ldots
\end{align}

\begin{center}
\adjustbox{max width=\textwidth}{%
\begin{tabular}{lc}
    \toprule
    & \textbf{Value} \\
    \midrule
    Measured: $\Delta n(1064 \to 2128\,\text{nm})$ & $0.0336$ \\
    T0 expression: $\sqrt{3}\cdot\xipar^{1/2}$     & $0.0200$ \\
    \bottomrule
\end{tabular}%
}
\end{center}

\medskip
The assignment $\Delta n = \sqrt{3}\cdot\xipar^{1/2}$ does not hold, because the
expression lies $41\,\%$ below the measured value. The exact equality with $0.020$
is an algebraic identity; the measured value is known to only 2--3 digits.

\clearpage

% ==============================================================================
\section{Further Matches on the \texorpdfstring{$\xipar$}{xi} Ladder}
\label{sec:further}
% ==============================================================================

\subsection{Ratio of Electro-Optic Coefficients}

LiNbO$_3$ has two relevant coupling constants: the linear electro-optic
(Pockels) coefficient $r_{33} = 30.8\,\text{pm/V}$ and the nonlinear
coefficient $d_{33} = 27.2\,\text{pm/V}$.  Their ratio:

\begin{equation}
    \frac{r_{33}}{d_{33}} = \frac{30.8}{27.2} = 1.1324
    \;\approx\; 1 + \xipar^{1/4} = 1.1075
    \qquad (\text{dev.}\ {+}2.2\,\%)
\end{equation}

$\xipar^{1/4}$ is the sub-Planck geometry exponent in T0 (Doc.~009).
The ratio of two electro-optic coupling constants lies on the $\xipar$
ladder — indicating that the crystal band structure itself is
$\xipar$-rational.

\clearpage

% ==============================================================================
\section{\texorpdfstring{$\xipar^{1/2}$}{xi\^{}(1/2)} as the Electroweak Rung}
\label{sec:classification}
% ==============================================================================

\begin{table}[H]
    \centering
    \adjustbox{max width=\textwidth}{%
\begin{tabular}{llll}
        \toprule
        \textbf{Domain} & \textbf{Quantity} &
        \textbf{T0 expression} & \textbf{Doc.} \\
        \midrule
        \multicolumn{4}{l}{\itshape Particle physics (Standard Model)} \\[2pt]
        & Weak coupling       & $\alpha_W = \xipar^{1/2}$
            & 041, 061 \\
        & Lepton masses       & $m_e/m_\mu \propto \xipar^{1/2}$
            & 056, 122, 133 \\
        & Bottom quark        & $m_b = v\cdot\tfrac{3}{2}\cdot\xipar^{1/2}$
            & 041 \\
        \midrule
        \multicolumn{4}{l}{\itshape Material properties (this document)} \\[2pt]
        & EO coefficients
            & $r_{33}/d_{33} \approx 1 + \xipar^{1/4}$
            & \textbf{167} \\
        \midrule
        \multicolumn{4}{l}{\itshape Cosmology and quantum optics} \\[2pt]
        & Casimir scale       & $L_\xi = (15\xipar/\pi^2)^{1/4}\,\hbar c/(k_BT_\text{CMB})$
            & 078, 166 \\
        & $H_0$ ratio         & $\hbar H_0/m_ec^2 = (\pi/2)\cdot\xipar^{10}$
            & 165 \\
        \bottomrule
    \end{tabular}%
}
    \caption{$\xipar^{1/2}$ as the electroweak exponent and further $\xipar$
    powers of T0 theory. The PPLN dispersion does not lie on the
    $\xipar^{1/2}$ rung (Section~\ref{sec:main_result}).}
    \label{tab:universal}
\end{table}

The PPLN dispersion does not belong to this hierarchy: $\sqrt{3}\cdot\xipar^{1/2}$
lies $41\,\%$ below the measured $\Delta n = 0.0336$. Likewise, the relations
$E_\mathrm{EW} \sim E_P\cdot\xipar^{1/2}$ and $m_h/M_P = \xipar^{1/2}$ are absent,
because they are not tenable (Section~\ref{sec:xi_half}).

\clearpage

% ==============================================================================
\section{Consequences for CIM Design}
\label{sec:consequences}
% ==============================================================================

A geometric answer to ``Why LiNbO$_3$?'' via $\Delta n = \sqrt{3}\cdot\xipar^{1/2}$
does not hold, because the measured dispersion is $0.0336$. The choice of material
remains empirically founded: largest known $d_{33}$ coefficient combined with low
absorption losses. A selection rule $\Delta n = k\cdot\xipar^{n/4}$ for CIM substrates
based on this assignment, and predictions derived from it
($\sqrt{2}\cdot\xipar^{1/2}$ at $\sim 900\,\text{nm}$, $2\cdot\xipar^{1/2}$ at
$\sim 680\,\text{nm}$), therefore do not hold.

\clearpage

% ==============================================================================
\section{Summary}
\label{sec:summary}
% ==============================================================================

\begin{enumerate}

    \item \textbf{Result}: The dispersion of LiNbO$_3$ for the most important
    PPLN process ($1064 \to 2128\,\text{nm}$) is $\Delta n = 0.0336$ according
    to Sellmeier (Zelmon 1997). The assignment
    $\Delta n = \sqrt{3}\cdot\xipar^{1/2} = 0.020$ does not hold: the expression
    lies $41\,\%$ below the measured value, and the equality with $0.020$ is an
    identity.

    \item \textbf{Exponent}: $\xipar^{1/2}$ appears as
    $\alpha_W = \xipar^{1/2}$, $m_e/m_\mu \propto \xipar^{1/2}$ and in further
    quantities (Docs.~041, 056, 061, 122, 133). $E_\mathrm{EW} \sim E_P\cdot\xipar^{1/2}$
    and $m_h/M_P = \xipar^{1/2}$ are not tenable. With the Casimir scale, which
    carries the exponent $1/4$ in Doc.~166, $\xipar^{1/2}$ shares no exponent.

    \item \textbf{Further match}: $r_{33}/d_{33} \approx 1 + \xipar^{1/4}$
    (dev.\ $2.2\,\%$).

    \item \textbf{Classification and conclusion}: The PPLN dispersion does not
    lie on the $\xipar^{1/2}$ rung. The choice of material for CIM design remains
    empirically founded.

\end{enumerate}

\bigskip\bigskip

\noindent\textbf{Related documents:}
Doc.~041 (Parameter derivation, $\alpha_W = \xipar^{1/2}$),
Doc.~056 (Universal derivation, lepton masses),
Doc.~061 (CMB units, $\alpha_W$),
Doc.~081 (Summary, electroweak scale),
Doc.~122 (Ratio vs.\ absolute, $m_e/m_\mu$),
Doc.~133 (Fractal correction),
Doc.~161 (CIM fundamentals),
Doc.~165 ($H_0$ ratio),
Doc.~166 (Casimir-CMB-$H_0$ bridge).
